% Part: incompleteness
% Chapter: theories-computability
% Section: computably-axiomatizable

\documentclass[../../../include/open-logic-section]{subfiles}

\begin{document}

\olfileid{inc}{tcp}{cax}

\olsection{\printtoken{S}{axiomatizable} نظریے}

نظریہ $\Th{T}$، \emph{!!{axiomatizable}} اے جے اوہدے
مسلّمیاں دا کوئی قابلِ حساب سیٹ~$A$ ہووے۔ ($A$ نوں $\Th{T}$ لئی
مسلّمیاں دا سیٹ آکھن دا مطلب اے کہ $T = \Setabs{!A}{A \Proves !A}$۔)
قدرتی عدداں دی کوئی وی “معقول” مسلّمہ بندی ایہ شرط پوری کرے گی۔
خاص طور اُتے ہر اوہ نظریہ !!{axiomatizable} اے جیہدے مسلّمیاں
دا سیٹ محدود اے۔

\begin{lem}
فرض کرو $\Th{T}$ !!{axiomatizable} اے۔ تاں $\Th{T}$
!!{computably
enumerable} اے۔
\end{lem}

\begin{proof}
فرض کرو $A$، $\Th{T}$ لئی مسلّمیاں دا قابلِ حساب سیٹ اے۔
$!A \in T$ دی تصدیق لئی مسلّمیاں توں $!A$ دا !!a{derivation}
لبھو۔ ثبوت مل جائے تاں رکنیت ثابت ہو جاندی اے۔ جے ثبوت موجود
نہ ہووے، تاں ایہ تلاش ضروری نہیں کہ مُک جائے۔

دوجے لفظاں وچ $!A$، $\Th{T}$ وچ اُدوں تے صرف اُدوں اے
جدوں مسلّمیاں دی اک محدود فہرست $!B_1$، \dots، $!B_k$
موجود ہووے، جیہدے سارے مسلّمے سیٹ $A$ وچ ہون، تے پہلے درجے
دی منطق وچ $(!B_1 \land \dots \land !B_k) \lif !A$ دا
!!a{derivation} ہووے۔ اسیں پہلاں جان دے آں کہ “کوئی \dots
موجود اے جیہدے لئی \dots” دے روپ وچ تعریف ہو سکّن والا ہر
سیٹ !!{c.e.} اے، جے دوجے “\dots” والی شرط قابلِ حساب ہووے۔
\end{proof}

\end{document}
